\documentclass{article}

\usepackage{amsmath}
\usepackage{graphicx}
\usepackage{hyperref}
\usepackage{booktabs}

\title{CSCI 5942: AI Engineering Assignment 2}
\author{Bill Hsu}
\date{\today}

\begin{document}

\maketitle

\section{Part 1: Wire in your base checkpoint}

Base model, nats/char on the four corpora (400 non-overlapping 256-char
blocks each):

\begin{center}
\begin{tabular}{lr}
\toprule
Corpus & nats/char \\
\midrule
shakespeare & 1.4321 \\
tinystories & 2.1625 \\
wikipedia   & 2.4880 \\
python-code & 4.8962 \\
\bottomrule
\end{tabular}
\end{center}

Sanity check: shakespeare (1.4321) lands within a few hundredths of the
A1 base training val loss (1.4668), as expected.

\begin{figure}[ht]
    \centering
    % TODO: replace with your screenshot of the Opik traces view (50 generations)
    \includegraphics[width=0.9\textwidth]{screenshots/part1_traces.png}
    \caption{Opik traces view showing the 50 base-model generations.}
    \label{fig:part1-traces}
\end{figure}

\section{Part 2: The grid}

\begin{figure}[ht]
    \centering
    \includegraphics[width=0.9\textwidth]{results/heatmap.png}
    \caption{nats/char, 5 model sizes (pico--base) $\times$ 4 corpora.}
    \label{fig:heatmap}
\end{figure}

OOV fractions (from \texttt{corpora/README.md}):

\begin{center}
\begin{tabular}{lr}
\toprule
Corpus & OOV deleted \\
\midrule
shakespeare.filtered.txt   & 0.00\% \\
tinystories.filtered.txt   & 0.98\% \\
wikipedia.filtered.txt     & 2.30\% \\
python-code.filtered.txt   & 7.69\% \\
\bottomrule
\end{tabular}
\end{center}

% TODO: revise these two sentences to match what your own heatmap shows
Only the shakespeare column improves monotonically with model size; the
python-code column is the worst regression, with our small pico model
outscoring several larger rungs. A bigger Shakespeare-trained model
learns more about Shakespeare's specific character distribution, but
that specialization does not transfer to a domain (Python source, with
its underscores, parentheses, and digits deleted as OOV) whose
character statistics look nothing like Elizabethan verse.

\section{Part 3: The twin}

\begin{figure}[ht]
    \centering
    \includegraphics[width=0.9\textwidth]{results/twin_curves.png}
    \caption{Twin (base config trained on TinyStories) train/val curves.}
    \label{fig:twin-curves}
\end{figure}

\begin{figure}[ht]
    \centering
    \includegraphics[width=0.9\textwidth]{results/heatmap_twin.png}
    \caption{Part 2 heatmap with the twin's row appended.}
    \label{fig:heatmap-twin}
\end{figure}

Predicted vs.\ measured:

\begin{center}
\begin{tabular}{lrrr}
\toprule
Corpus & Predicted & Measured & Error \\
\midrule
tinystories & 1.46 & 0.6868 & 0.77 \\
shakespeare & 2.58 & 2.6345 & 0.05 \\
\bottomrule
\end{tabular}
\end{center}

% TODO: revise to match your own reasoning
The larger miss was on tinystories: TinyStories' simple, highly
repetitive sentence structure makes it far more predictable at the
character level than the prediction assumed, and that low intrinsic
entropy outweighs the disadvantage of the twin seeing each character
only once (vs.\ the base model's roughly 74 epochs over Shakespeare).

\section{Part 4: The judge}

Pointwise means with 95\% CIs, $n=50$:

\begin{center}
\begin{tabular}{llrrr}
\toprule
Model & Rubric & grammar & coherence & style\_match \\
\midrule
base             & shakespeare & 2.24 $\pm$ 0.143 & 1.20 $\pm$ 0.112 & 4.30 $\pm$ 0.405 \\
base             & tinystories & 1.56 $\pm$ 0.139 & 1.00 $\pm$ 0.000 & 1.00 $\pm$ 0.000 \\
base-tinystories & shakespeare & 5.02 $\pm$ 0.477 & 2.84 $\pm$ 0.523 & 1.00 $\pm$ 0.000 \\
base-tinystories & tinystories & 5.08 $\pm$ 0.454 & 2.80 $\pm$ 0.336 & 4.70 $\pm$ 0.469 \\
\bottomrule
\end{tabular}
\end{center}

Pairwise win rates (base vs.\ base-tinystories), with 95\% CIs and order-swap flip counts, $n=50$ pairs each:

\begin{center}
\begin{tabular}{lrrr}
\toprule
Rubric & base win rate & base-tinystories win rate & flips \\
\midrule
shakespeare & 1.00 $\pm$ 0.00 (n=49) & 0.00 $\pm$ 0.00 (n=49) & 1 \\
tinystories & 0.00 $\pm$ 0.00 (n=50) & 1.00 $\pm$ 0.00 (n=50) & 0 \\
\bottomrule
\end{tabular}
\end{center}

\begin{figure}[ht]
    \centering
    % TODO: replace with your screenshot of the scored traces
    \includegraphics[width=0.9\textwidth]{screenshots/part4_scored_traces.png}
    \caption{Opik traces view with judge feedback scores attached.}
    \label{fig:part4-scored-traces}
\end{figure}

\section{Part 5: Perturbing the judge}

All three runs score the same 50 base-model samples ($n=50$); only the judge prompt changes.

\begin{center}
\begin{tabular}{llrrr}
\toprule
Prompt version & Rubric & grammar & coherence & style\_match \\
\midrule
original (Part 4)  & shakespeare & 2.24 $\pm$ 0.143 & 1.20 $\pm$ 0.112 & 4.30 $\pm$ 0.405 \\
paraphrase          & shakespeare\_v2 & 2.56 $\pm$ 0.187 & 1.38 $\pm$ 0.147 & 4.42 $\pm$ 0.400 \\
axes-reordered       & shakespeare & 1.80 $\pm$ 0.186 & 1.40 $\pm$ 0.137 & 4.32 $\pm$ 0.377 \\
\bottomrule
\end{tabular}
\end{center}

% TODO: revise to match your own framing
style\_match and coherence stay within overlapping confidence intervals
across all three prompt versions, so the Part 4 quality comparison
(which rests on style\_match) survives both perturbations. grammar does
not: its confidence interval under axes-reordered (1.61--1.99) does not
overlap the original (2.10--2.38), so simply moving grammar from first
to last in the axis list significantly lowers its score even though the
axis description itself never changed -- any Part 4 claim resting on the
grammar number specifically should be treated as within the judge's own
prompt-order noise.

\section{The Question}

% TODO: half a page max. Cite the Part 2 grid for the loss claim and the
% Part 4 judge tables for the quality claim, and point at where they disagree.
TODO: which model is better, and on what exam?

\end{document}
